\section{Forward and Backward Propagation}
\label{sec:propagation}

We now endow the directed simplex memory graph with hidden states and propagation laws.

\subsection{General forward propagation}

Let each hidden node $\alpha\in V_{n,m}$ carry a state $h_\alpha\in\R^d$. Let the external input be $u\in\R^{d_{\mathrm{in}}}$ and the external output be $y\in\R^{d_{\mathrm{out}}}$. We denote the hidden activation by $\sigma$ and the readout activation by $\phi$.

A natural forward propagation law on the directed hidden graph is
\begin{equation}\label{eq:general-forward}
 h_\alpha
 =
 \sigma\!\left(
 \sum_{\gamma\to\alpha} W_{\alpha\gamma} h_\gamma
 +
 (B_{\mathrm{in}}u)_\alpha\,\1_{\{\alpha\in F_{\mathrm{in}}\}}
 +
 b_\alpha
 \right),
 \qquad \alpha\in V_{n,m}.
\end{equation}
The readout is taken from the output facet:
\begin{equation}\label{eq:general-readout}
 y
 =
 \phi\!\left(
 C_{\mathrm{out}}\, h|_{F_{\mathrm{out}}}
 +
 b_{\mathrm{out}}
 \right).
\end{equation}
Here $h|_{F_{\mathrm{out}}}$ denotes the concatenation of hidden states on the output facet.

\subsection{Triangle and tetrahedron forward laws}

The minimal motifs can be written explicitly.

\begin{example}[Triangle forward law]
For the triangle motif, let $h_a,h_c,h_b\in\R^d$ denote the hidden states. A minimal forward law is
\begin{align}
 h_a &= \sigma(W_a u + b_a), \label{eq:triangle-a}\\
 h_c &= \sigma(W_{ca} h_a + U_c u + b_c), \label{eq:triangle-c}\\
 h_b &= \sigma(W_{ba} h_a + W_{bc} h_c + b_b), \label{eq:triangle-b}\\
 y &= \phi(R_b h_b + R_c h_c + b_{\mathrm{out}}). \label{eq:triangle-y}
\end{align}
The term $R_c h_c$ expresses direct readout from the shared interface, while the term $W_{ba}h_a$ expresses a skip path along the backbone.
\end{example}

\begin{example}[Tetrahedron forward law]
For the tetrahedron motif, let $h_a,h_b\in\R^d$ and let $h_C\in\R^{2d}$ denote the concatenated state of the two-node intermediate layer $\{c,d\}$. A minimal forward law is
\begin{align}
 h_a &= \sigma(W_a u + b_a), \label{eq:tetra-a}\\
 h_C &= \sigma(W_{Ca} h_a + U_C u + b_C), \label{eq:tetra-C}\\
 h_b &= \sigma(W_{ba} h_a + W_{bC} h_C + b_b), \label{eq:tetra-b}\\
 y &= \phi(R_b h_b + R_C h_C + b_{\mathrm{out}}). \label{eq:tetra-y}
\end{align}
Again, $h_C$ admits direct readout, while $h_b$ represents the more integrated terminal hidden state on the backbone.
\end{example}

\subsection{Backward propagation and bidirectional coupling}

Let $L=L(y)$ be a loss function. For the tetrahedron motif, write
\[
\delta_b := \frac{\partial L}{\partial h_b},
\qquad
\delta_C := \frac{\partial L}{\partial h_C},
\qquad
\delta_a := \frac{\partial L}{\partial h_a}.
\]
Then the backpropagated signals satisfy relations of the form
\begin{align}
 \delta_C
 &=
 R_C^\top \frac{\partial L}{\partial y}
 +
 W_{bC}^\top \delta_b,
 \label{eq:backprop-C}
 \\
 \delta_a
 &=
 W_{Ca}^\top \delta_C
 +
 W_{ba}^\top \delta_b.
 \label{eq:backprop-a}
\end{align}
Thus the intermediate short-term layer receives training signals both directly from the readout and indirectly through the terminal backbone state, while the initial backbone node is shaped by both branches. In this sense, the architecture encodes a two-way interaction between fast and slow memory pathways already at the level of backpropagation.
